\scanpage{213}\mspage{158}

% MS 1–5
est licite) -- \deleted{\unsure{alors}} la description de l'image inv. est OK pour \uline{\unsure{tt}} chgt de base. Mais on a vu que dans le cas d'un \unsure{voisin|voisinage} local comme $ {}_{s\backslash}S'\longrightarrow S'$, il n'y a pas même besoin de cet artifice -- la formule est OK sans restriction de lissité.

% MS 6
\deleted{\unsure{Comme le}}\inserted{Ainsi}

% MS 7
\[
 q^*(\unsure{\xi|\zeta}')\simeq\mathrm{hot}_{{}_{s\backslash}S',\unsure{W|w}}({}_{s\backslash}X'),
\]

% MS 8
d'où

% MS 9
\[
 f'_!(q^*(\xi'))\simeq\mathrm{hot}_{\unsure{e|S},W}({}_{s\backslash}X')=\mathrm{hot}_W({}_{s\backslash}X').
\]

% MS 10
Donc on a bien trouvé

% MS 11
\[
 p^*f_!\simeq f'_{\unsure{!|!*}}q^*.
\]

% MS 12–15
On peut maintenant prouver aussi Der 3 b. On a vu qu' \deleted{\unsure{une}} fois qu'on dispose des adjoints à gauche $f_!$ des $f^*$, cet axiome s'énonce sous \uline{forme homologique} ainsi :

% MS 16
Si

% MS 17
\[
 f:S'\longrightarrow S
\]

% MS 18–19
est dans $W$, il faut prouver que pour tout

% MS 19
\[
 \xi\in\mathrm{HOT}_W(e)\overset{\mathrm{def}}{=}\mathrm{Hot}_W
\]

% MS 20
la flèche correspondante

% MS 21
\[
 \unsure{f'|p'}_!\,p'^*(\xi)\longrightarrow p_!p^*(\xi)\tag*{$(*)$}
\]

% MS 22–25
est iso dans $\mathrm{Hot}_W$, où $p:S\longrightarrow e$, $p':S'\longrightarrow e$ sont les flèches structurales. Or comme tout objet de $\mathrm{Cat}$ est lisse sur $e$, le cor.~2. précédent permet d'identifier cette flèche \unsure{bis|bas|les} :

% MS 26
\[
 \unsure{s|S}'\times_S\xi\longrightarrow s\times_S\xi
\]

% archive sheet number: 212
